\documentclass[journal]{IEEEtran}
\usepackage{amsmath}
\usepackage{amssymb}
\usepackage{graphicx}
% -----------------------------------------------------------------------
% Packages
% -----------------------------------------------------------------------
\usepackage{booktabs}
\usepackage{array}
\usepackage{amsmath,amsfonts,amssymb}
\usepackage{float}
\usepackage{array}
\usepackage[caption=false,font=normalsize,labelfont=sf,textfont=sf]{subfig}
\usepackage{textcomp}
\usepackage{stfloats}
\usepackage{url}
\usepackage{verbatim}
\usepackage{graphicx}
\usepackage{cite}
\usepackage{booktabs}
\usepackage{tabularx}
\usepackage{multirow}
\usepackage{ragged2e}
\usepackage{float}
\usepackage{placeins}        % \FloatBarrier
\usepackage[caption=false,font=normalsize,labelfont=sf,textfont=sf]{subfig}

%% -------- Color (performance-table shading) ----------------
\usepackage[table]{xcolor}
\usepackage{colortbl}

\definecolor{darkred}{RGB}{153,0,0}
\definecolor{lightred}{RGB}{255,230,230}
\definecolor{midred}{RGB}{255,165,165}
\definecolor{headerblue}{RGB}{30,80,140}
\definecolor{rowgray}{RGB}{240,240,240}

%% -------- Custom column type -------------------------------
\newcolumntype{Y}{>{\RaggedRight\arraybackslash}X}

%% -------- Image path (all figures live in images/) ---------
\graphicspath{{images/}}

\hyphenation{op-tical net-works semi-conduc-tor IEEE-Xplore Wire-less}

% Custom algorithm float (no external algorithm package required)
\floatstyle{ruled}
\newfloat{algorithm}{tbp}{loa}
\floatname{algorithm}{Algorithm}

% -----------------------------------------------------------------------
\begin{document}
% -----------------------------------------------------------------------

\title{DDQN for Coverage and Battery Optimization\\in WSN Scheduling}

\author{Vinay~Kumar,~Siddarth~Vaidya,~Ch~Tarun~Kumar,~Ankitesh~Shrivastava%
}

% Running headers
\markboth{IEEE Transactions on Wireless Communications}%
{Kumar \MakeLowercase{\textit{et al.}}: DDQN for Coverage and Battery
Optimization in WSN Scheduling}



\maketitle

% -----------------------------------------------------------------------
\begin{abstract}
Wireless Sensor Networks (WSNs) are widely deployed for environmental
monitoring, industrial automation, and smart infrastructure, yet their
operational lifetime remains fundamentally constrained by battery energy
and irreversible cell degradation. Classical scheduling
policies—including greedy wake-up, round-robin, and static threshold
approaches—fail to account for the nonlinear coupling between energy
consumption, battery State of Health (SoH), and spatial coverage.
Standard Deep Q-Network (DQN) approaches partially address this via
learned scheduling but are theoretically susceptible to systematic
Q-value overestimation, a property whose impact is expected to grow
with the dimensionality and variance of multi-node, multi-objective
settings.

This paper presents \textbf{DDQN-WSN}, a constraint-aware
\textit{factorized} Double Deep Q-Network agent trained in a Gymnasium
simulation of a 50-node WSN deployed over a
$500\,\text{m} \times 500\,\text{m}$ arena and equipped with an empirical
SoC/SoH-aware battery degradation model. The system model incorporates
per-node battery dynamics (cycle-based and calendar SoH degradation), a
charging state machine, and a cooperative wake-up rule between
energy-depleted nodes. To avoid the combinatorial intractability of the
full $2^N$ joint action space, scheduling is formulated as $N$
parameter-shared per-node binary decisions; the resulting raw action
vector is passed through a constraint-aware execution layer that enforces
charging feasibility and coverage continuity before the physical state
transition. The agent observes a 9-feature per-node state vector—including
normalized node coordinates and a local awake-neighbour density score
that exposes spatial geometry to the policy—and outputs per-node
SLEEP/AWAKE decisions. Reward is a weighted sum of four objectives:
area coverage, SoC-weighted energy efficiency, battery health
preservation, and charge-level fairness. Ablation against standard DQN under identical environment and
hyperparameter conditions shows that the two agents reach
near-parity on every converged metric measured at training time:
both exceed $96\%$ mean area coverage and neither breaches the $30\%$
node-death threshold within an episode across the entire 500-episode
training run. DDQN-WSN is therefore reported as achieving performance
comparable to DQN under the present setting, while providing a more
robust target-estimation mechanism that decouples action selection
from action evaluation and is theoretically less susceptible to the
systematic Q-value overestimation that compounds across $N$
parallel per-node argmax operations---a stability property that is
expected to matter most in longer-horizon or higher-variance
deployment settings beyond the scope of this study.
\end{abstract}

\begin{IEEEkeywords}
Wireless Sensor Networks, Double Deep Q-Network, Sleep Scheduling,
Battery Degradation, State of Health, Reinforcement Learning, Network
Lifetime Optimization.
\end{IEEEkeywords}

% -----------------------------------------------------------------------
\section{Introduction}
\label{sec:intro}
% -----------------------------------------------------------------------

\IEEEPARstart{W}{ireless} Sensor networks deploy large numbers of
battery-powered nodes to sense and relay information from physical
environments. Applications range from agricultural monitoring and
structural health sensing to industrial IoT and military surveillance.
Despite decades of research, the fundamental tension between sensing duty
and battery longevity remains unsolved at scale: every step a node is
active drains its battery and accelerates electrochemical degradation,
yet keeping too many nodes asleep degrades spatial coverage and data
continuity.
\IEEEpubidadjcol
Energy depletion in WSNs is governed by two coupled processes: the
\textit{State of Charge} (SoC), which represents immediately available
energy and can be partially recovered through charging, and the
\textit{State of Health} (SoH), which captures cumulative cell
degradation from deep discharge cycles and calendar fade and is strictly
irreversible. Existing scheduling algorithms typically optimize for one
objective in isolation—greedy policies maximize instantaneous coverage
by keeping the maximum number of nodes awake, rapidly exhausting
batteries and causing premature SoH collapse; round-robin policies
achieve fairness but ignore spatial redundancy, wasting energy in
already-covered regions; and threshold-based heuristics react to SoC
levels but cannot anticipate future demand, yielding suboptimal
trade-offs between network lifetime and coverage. None of these policies
model battery degradation physics, cooperative inter-node behavior, or
the long-horizon consequences of scheduling decisions. Deep Q-Networks
offer a principled approach to sequential scheduling via function
approximation over state-action value space; however, standard DQN uses
the same network to both \textit{select} and \textit{evaluate} the
next-state action, producing systematic Q-value overestimation that is
particularly pronounced in high-dimensional, correlated per-node action
spaces. This overestimation is a well-documented theoretical
failure mode of standard DQN whose practical impact is expected to
grow with the dimensionality and variance of the multi-node,
multi-objective WSN scheduling problem. To address these gaps, this
work contributes the following: (i)~an empirical SoC/SoH-aware
Gymnasium-compliant WSN simulation environment incorporating a surrogate
per-node battery degradation model—capturing the dominant SoC depletion,
depth-of-discharge (DoD) cycle aging, and calendar fade
mechanisms—together with a threshold-triggered charging state machine
and a cooperative wake-up rule that transfers sensing responsibility
between energy-depleted and healthy neighbors; (ii)~a compact
nine-feature centralized per-node observation vector encoding SoC, SoH,
last action, normalized distance to sink, normalized $(x, y)$
coordinates, a local awake-neighbour density score, exponential moving
average activity ratio, and charging flag, exposing both energy state
and spatial geometry directly to the policy; (iii)~a constraint-aware
factorized action formulation in which the agent's raw per-node binary
output $\mathbf{a}_t$ is filtered by a deterministic execution layer
into the executed action $\tilde{\mathbf{a}}_t = \mathcal{T}(s_t,
\mathbf{a}_t)$ that enforces charging feasibility and coverage
continuity, with replay transitions stored using $\tilde{\mathbf{a}}_t$;
(iv)~a weighted four-component reward function jointly optimizing area
coverage, energy efficiency, battery health, and charge-level fairness;
(v)~application of the Double DQN target to the factorized per-node
Q-head, decoupling action selection (online network) from value
evaluation (target network) so that the per-node maximization bias
does not compound additively across the $N$ parallel argmax
operations of the factorized head; and (vi)~a direct DDQN-vs-DQN
empirical ablation under identical environment and hyperparameter
conditions, characterising the parity of converged behaviour in the
500-episode setting and the theoretical robustness benefit of the
decoupled target on the factorized scheduling problem.

% -----------------------------------------------------------------------
\section{Related Work}
\label{sec:related}
% -----------------------------------------------------------------------

\subsection{Coverage- and Energy-Focused Scheduling}

Recent works primarily focus on optimizing energy consumption and
network lifetime using lightweight optimization and clustering-based
strategies. For instance, soft clustering approaches such as
\cite{burger2023} improve network lifetime but suffer from centroid
sensitivity, while fuzzy-logic-based routing methods
\cite{nguyen2024} enhance adaptability at the cost of increased
rule-processing overhead. Similarly, PSO-based hybrid approaches
\cite{chang2023} attempt to optimize routing decisions but remain
highly sensitive to parameter tuning.

Although these methods improve energy efficiency, they rely on heuristic
or rule-based decision-making and do not scale well with increasing
network complexity. Moreover, none of these approaches consider battery
State-of-Health (SoH), limiting their ability to handle long-term
energy degradation.

\subsection{Clustering and Scheduling Protocols}

Clustering and scheduling mechanisms remain widely adopted for reducing
communication overhead and balancing energy consumption. Modified LEACH
variants \cite{kaur2021} and duty-cycling strategies based on
reinforcement learning \cite{muhammad2023} aim to extend node lifetime
through adaptive activity scheduling. Additionally, packet scheduling
approaches such as QDNC \cite{li2023} optimize delay and delivery
performance.

However, these approaches exhibit several limitations, including slow
convergence, scalability challenges, and lack of adaptability to dynamic
network conditions. Most importantly, cluster formation and scheduling
decisions are typically based on residual energy or predefined rules,
without incorporating battery degradation effects such as SoH.

\subsection{Reinforcement Learning-Based Approaches}

Reinforcement learning techniques have been increasingly applied to WSN
routing and scheduling. Standard DQN-based approaches \cite{xue2020}
enable learning-based optimization of coverage and lifetime but suffer
from Q-value overestimation and limited generalization. Q-learning-based
adaptive routing methods \cite{su2023} provide improved path selection
but rely on tabular representations, making them unsuitable for
large-scale networks.

More recent approaches, such as federated learning combined with
clustering \cite{mohammed2024}, attempt to improve scalability and
distributed learning but introduce high computational complexity and
communication overhead.

Despite these advancements, existing RL-based solutions do not model
battery SoH and largely ignore long-term degradation effects. As a
result, their optimization remains short-sighted and suboptimal in
realistic deployments.

\subsection{Positioning of This Work}

The proposed DDQN-WSN framework addresses the key limitations identified
in prior work. Unlike existing approaches, it explicitly incorporates
battery State-of-Health (SoH) into the reward function, enabling
long-term energy-aware decision-making. Furthermore, it employs Double
DQN to eliminate Q-value overestimation observed in standard DQN-based
methods. In contrast to heuristic, fuzzy, and clustering-based
approaches, the proposed method learns adaptive policies directly from
the environment.

Additionally, the framework introduces a cooperative wake-up mechanism
that improves coverage and coordination without incurring additional
communication overhead. These features collectively overcome the
limitations of scalability, adaptability, and lack of battery awareness
observed in existing methods, as summarized in
Table~\ref{tab:related}.

\begin{table*}[!t]
\caption{Comparison of Related Works\label{tab:related}}
\centering
\renewcommand{\arraystretch}{1.25}
\small
\begin{tabular}{|p{2.5cm}|c|p{2.8cm}|p{2.2cm}|c|p{3.8cm}|}
\hline
\rowcolor{headerblue}
\textcolor{white}{\textbf{Work}} &
\textcolor{white}{\textbf{Year}} &
\textcolor{white}{\textbf{Method}} &
\textcolor{white}{\textbf{Key Metric}} &
\textcolor{white}{\textbf{SoH Model}} &
\textcolor{white}{\textbf{Limitation}} \\
\hline

Mohammed et al.~\cite{mohammed2024} & 2024 &
Federated learning + clustering &
Energy / overhead &
No &
High complexity; no battery health model \\
\hline

\rowcolor{rowgray}
Burger~\cite{burger2023} & 2023 &
Soft K-Means clustering &
Network lifetime &
No &
Centroid sensitivity \\
\hline

Nguyen et al.~\cite{nguyen2024} & 2024 &
Fuzzy logic cluster routing &
Lifetime / energy &
No &
Fuzzy-rule overhead \\
\hline

\rowcolor{rowgray}
Kaur et al.~\cite{kaur2021} & 2021 &
Modified LEACH &
Lifetime / stability &
No &
Dynamic-network control overhead \\
\hline

Xue Lin~\cite{xue2020} & 2020 &
Standard DQN &
Lifetime / coverage &
No &
Q-value overestimation; no SoH model \\
\hline

\rowcolor{rowgray}
Li et al.~\cite{li2023} & 2023 &
QDNC packet scheduling &
Delay / delivery &
No &
Scalability not demonstrated \\
\hline

Muhammad et al.~\cite{muhammad2023} & 2023 &
RL sleep/wake duty cycling &
Node lifetime &
No &
Slow convergence \\
\hline

\rowcolor{rowgray}
Chang and Lee~\cite{chang2023} & 2023 &
PSO + fuzzy clustering &
Energy / routing &
No &
Parameter sensitivity \\
\hline

Su et al.~\cite{su2023} & 2023 &
Q-Learning adaptive routing &
Lifetime / paths &
No &
Tabular; not scalable \\
\hline

\textbf{DDQN-WSN (Proposed)} & \textbf{---} &
\textbf{Double DQN + SoH reward} &
\textbf{Lifetime / SoH} &
\textbf{Yes} &
Centralized; requires global state \\
\hline

\end{tabular}
\end{table*}

% -----------------------------------------------------------------------
\section{System Model}
\label{sec:system}
% -----------------------------------------------------------------------

\subsection{Network Model}
\label{sec:network_model}

We consider a WSN comprising $N = 50$ sensor nodes deployed uniformly
at random over a two-dimensional arena of size
$500\,\text{m} \times 500\,\text{m}$. A single sink node is fixed at
the arena centre $(250, 250)$. Node positions are drawn from a uniform
distribution at the start of each training episode and remain static
thereafter.

Each node $i$ occupies position $\mathbf{p}_i \in \mathbb{R}^2$. Its
normalized distance to the sink is

\begin{equation}
  d_i = \frac{\|\mathbf{p}_i - \mathbf{p}_{\text{sink}}\|_2}
             {\sqrt{W^2 + H^2}},
  \label{eq:dist}
\end{equation}

where $W = H = 500\,\text{m}$. This normalized distance is used in
both the observation vector and the distance-weighted energy draw
model.

At each discrete timestep $t$, the scheduling agent assigns a binary
action $a_i^t \in \{0, 1\}$ to every node, where $0$ denotes SLEEP and
$1$ denotes AWAKE. The episode terminates either when the number of dead
nodes exceeds the threshold $\delta \cdot N$ (with $\delta = 0.3$), or when
$t = T_{\max} = 1000$ steps.

Spatial coverage is evaluated against a $20 \times 20$ uniform grid of
sample points across the arena. A grid point $\mathbf{g}$ is considered
covered if at least one alive AWAKE node lies within the sensing radius
$r_s = 100\,\text{m}$:

\begin{equation}
  \text{coverage}(t) = \frac{1}{|\mathcal{G}|}
  \sum_{\mathbf{g} \in \mathcal{G}}
  \mathbf{1}\!\Big[\exists\, i :
  a_i^t {=} 1,\,
  \neg\text{dead}(i),\,
  \|\mathbf{p}_i {-} \mathbf{g}\|_2 {\leq} r_s \Big].
  \label{eq:coverage}
\end{equation}

Besides, in order to explicitly describe the overlapping of the sensing area and the efficiency of the space coverage, the Coverage Redundancy Metric (CRM) is introduced which is defined as the mean number of active sensing nodes that covers one of the covered grid points.

\begin{equation}
  \Gamma(t) =
  \begin{cases}
    \dfrac{\displaystyle\sum_{\mathbf{g} \in \mathcal{G}}
           \sum_{i=1}^{N} \mathbf{1}\big[ a_i^t = 1,\,
           \neg\text{dead}(i),\,
           \|\mathbf{p}_i - \mathbf{g}\|_2 \leq r_s \big]}
          {\displaystyle\sum_{\mathbf{g} \in \mathcal{G}}
           \mathbf{1}\big[\exists j : a_j^t = 1,\,
           \neg\text{dead}(j),\,
           \|\mathbf{p}_j - \mathbf{g}\|_2 \leq r_s \big]},
      & \text{if } \text{coverage}(t) > 0, \\[6pt]
    0, & \text{otherwise.}
  \end{cases}
  \label{eq:redundancy}
\end{equation}

Although the agent directly optimises the raw coverage, by monitoring the coverage metric, one can ensure that the learnt policy has a low spatial overlap (near 1 for the redundancy) without relying on brute redundancy in order to achieve continuous spatial coverage.

\subsection{Battery Model}

The networks considered in this paper are \emph{rechargeable WSNs}:
each node is equipped with an on-board energy buffer and may enter a
low-power charging state that recovers SoC from an assumed local
replenishment source (e.g.\ a solar harvester, a periodic
wireless-power transmitter, or a mobile charger). The replenishment
source itself is abstracted---only its effect on per-node SoC is
modelled---so that the same scheduler applies regardless of the
specific harvesting hardware. The full charging state machine is
described in Sec.~\ref{sec:exec_layer}.

Each node $i$ carries an independent \texttt{BatteryModel} instance
initialized at full charge ($\text{SoC}_i = E_{\max} = 100$) and full
health ($\text{SoH}_i = 1.0$).

\subsubsection{State of Charge (SoC)}

SoC represents immediately available energy. Energy draw per timestep
is action- and distance-dependent:

\begin{equation}
  e_i^t =
  \begin{cases}
    e_{\text{awake}} \cdot (1 + 0.1 \cdot d_i) &
      \text{if } a_i^t = 1\ (\text{AWAKE}), \\
    e_{\text{sleep}} &
      \text{if } a_i^t = 0\ (\text{SLEEP}),
  \end{cases}
  \label{eq:energy_draw}
\end{equation}

with $e_{\text{awake}} = 1.0$ and $e_{\text{sleep}} = 0.01$. The
distance penalty models the higher transmit power required by nodes far
from the sink. SoC is updated as:

\begin{equation}
  \text{SoC}_i^{t+1} = \max\!\left(0,\;
  \text{SoC}_i^t - e_i^t\right).
  \label{eq:soc_update}
\end{equation}

\subsubsection{State of Health: Cycle-Based Degradation}

Cycle-based State of Health (SoH) loss is modelled as a super-linear
function of the Depth of Discharge (DoD) traversed at each timestep,
capturing the empirically observed asymmetry whereby deep discharge
cycles damage cells disproportionately more than shallow ones. All
nodes are assumed to operate under nominal isothermal conditions with
homogeneous cell chemistry and a standard low-rate discharge profile;
the cycle-aging law used in the simulation therefore reduces to a
single DoD-dependent term parameterised by a baseline rate
$k_{\text{cycle}}$ and an exponent $\alpha > 1$:

\begin{equation}
  \Delta\text{SoH}_i^{\text{cycle}} = k_{\text{cycle}} \cdot
  \left( \frac{|\text{SoC}_i^t - \text{SoC}_i^{t+1}|}{E_{\max}}
  \right)^{\!\alpha}.
\label{eq:soh_cycle}
\end{equation}

Extensions that incorporate cell chemistry, Arrhenius-type
temperature dependence, or C-rate corrections are straightforward
in principle but are out of scope for the present surrogate
simulation; the parameter values used in the experiments are listed
in Table~\ref{tab:config}.

\subsubsection{SoH: Calendar Aging}

In addition to cycle-induced damage, SoH decreases passively at each
timestep due to electrolyte decomposition and self-discharge
reactions. Under the same controlled reference conditions, the
calendar-aging law reduces to a constant per-step fade rate
$k_{\text{cal}}$:

\begin{equation}
  \Delta\text{SoH}_i^{\text{cal}} = k_{\text{cal}}.
  \label{eq:soh_cal}
\end{equation}

\subsubsection{Combined SoH Update and Node-Death Criterion}

The combined SoH update per timestep is:

\begin{equation}
  \text{SoH}_i^{t+1} = \max\!\left(0,\;
  \text{SoH}_i^t
  - \Delta\text{SoH}_i^{\text{cycle}}
  - \Delta\text{SoH}_i^{\text{cal}}\right).
  \label{eq:soh_update}
\end{equation}

SoH is strictly non-increasing and does not recover during charging. A
node $i$ is declared \emph{dead} when either

\begin{equation}
  \text{SoC}_i \leq 0.01 \quad \text{or} \quad \text{SoH}_i \leq 0.05.
  \label{eq:death}
\end{equation}

Once dead, the node is permanently excluded from scheduling and coverage
calculations.

% -----------------------------------------------------------------------
\section{Proposed Methodology: Constraint-Aware Factorized DDQN for WSN Scheduling}
\label{sec:ddqn}
\begin{figure}[h]
    \centering
    \includegraphics[width=0.5\textwidth]{Flow.jpeg}
    \caption{End-to-end pipeline of the proposed
    constraint-aware factorized DDQN scheduler. The centralized
    per-node observation $\mathbf{s}_t$ concatenates nine features per
    node (SoC, SoH, previous executed action $\tilde{a}_i^{\,t-1}$,
    normalised distance to sink $d_i$, normalised coordinates
    $x_i/W$ and $y_i/H$, local awake-neighbour density
    $n_i^{\text{awake}}$, EMA activity ratio $\rho_i^t$, and charging
    flag $c_i^t$); the parameter-shared DDQN backbone emits the
    factorized raw action $\mathbf{a}_t \in \{0,1\}^N$ via per-node
    argmax; the constraint-aware execution layer
    $\mathcal{T}(\mathbf{s}_t, \mathbf{a}_t)$ produces the executed
    action $\tilde{\mathbf{a}}_t$ (charging-feasibility and
    coverage-continuity rules); the WSN environment applies
    $\tilde{\mathbf{a}}_t$, returning $(\mathbf{s}_{t+1}, r_t, d_t)$;
    and the tuple $(\mathbf{s}_t, \tilde{\mathbf{a}}_t, r_t,
    \mathbf{s}_{t+1}, d_t)$ is stored in the replay buffer that
    drives the Double DQN target update.}
    \label{fig:Flow diagram}
\end{figure}

The methodology is organized around five design choices that together
define the proposed agent: a centralized observation that exposes both
energy state and spatial geometry (Sec.~\ref{sec:state_rep}); a
factorized per-node binary action formulation that avoids the
combinatorial $2^N$ joint action space (Sec.~\ref{sec:action_factor}); a
deterministic constraint-aware execution layer that turns the agent's
raw action into a feasibility-respecting executed action used for both
transitions and replay (Sec.~\ref{sec:exec_layer}); a multi-objective
reward (Sec.~\ref{sec:reward}); and a parameter-shared DDQN learning
architecture with the standard Double DQN target adapted to the
factorized Q-head (Sec.~\ref{sec:ddqn_arch}).

% -----------------------------------------------------------------------

\subsection{Centralized State Representation}
\label{sec:state_rep}

At each timestep $t$, the environment returns a flat observation vector
$\mathbf{s}^t \in \mathbb{R}^{9N}$ obtained by concatenating a
9-dimensional feature block per node. The observation is centralized: a
single agent receives the full network state. For node $i$, the nine
features are listed in Table~\ref{tab:features}.

\begin{table}[!t]
\caption{Per-Node State Features\label{tab:features}}
\centering
\footnotesize
\begin{tabular}{clcp{2.4cm}}
\toprule
\textbf{Index} & \textbf{Feature} & \textbf{Range} &
\textbf{Description} \\
\midrule
$9i{+}0$ & $\text{SoC}_i / E_{\max}$    & $[0, 1]$
  & Norm.\ available energy \\
$9i{+}1$ & $\text{SoH}_i$               & $[0, 1]$
  & Battery health (monoton.\ decr.) \\
$9i{+}2$ & $\tilde{a}_i^{\,t-1}$        & $\{0, 1\}$
  & Prev.\ executed action \\
$9i{+}3$ & $d_i$                        & $[0, 1]$
  & Norm.\ dist.\ to sink \\
$9i{+}4$ & $x_i / W$                    & $[0, 1]$
  & Norm.\ x-coordinate \\
$9i{+}5$ & $y_i / H$                    & $[0, 1]$
  & Norm.\ y-coordinate \\
$9i{+}6$ & $n_i^{\text{awake}}$         & $[0, 1]$
  & Local awake-neighbour density \\
$9i{+}7$ & $\rho_i^t$                   & $[0, 1]$
  & EMA activity ratio \\
$9i{+}8$ & $c_i^t$                      & $\{0, 1\}$
  & Charging flag \\
\bottomrule
\end{tabular}
\end{table}

The activity ratio is updated per step as:

\begin{equation}
  \rho_i^{t+1} = 0.9 \cdot \rho_i^t + 0.1 \cdot \tilde{a}_i^{\,t},
  \label{eq:activity}
\end{equation}

providing the agent with a smoothed estimate of each node's recent duty
cycle. The local awake-neighbour density score is defined as the
fraction of node $i$'s in-radius neighbours that are alive and AWAKE
at step $t$. Let
\begin{equation}
  \mathcal{N}_i = \big\{ j \neq i :
  \|\mathbf{p}_i - \mathbf{p}_j\|_2 \leq r_s \big\}
  \label{eq:neighbourhood}
\end{equation}
denote the (static) set of geometric neighbours of node $i$ within
sensing radius $r_s$. The awake-neighbour density is then

\begin{equation}
  n_i^{\text{awake}} = \frac{\Big|\big\{ j \in \mathcal{N}_i :\,
        \tilde{a}_j^{\,t} = 1,\, \neg\text{dead}(j) \big\}\Big|}
        {|\mathcal{N}_i| + \varepsilon},
  \label{eq:awake_density}
\end{equation}

with $\varepsilon$ a small positive constant ($\varepsilon \ll 1$,
e.g.\ $10^{-6}$) that keeps the expression well-defined when a node
has no geometric neighbours within $r_s$ (in which case $n_i^{\text{awake}}
\to 0$). This normalisation gives $n_i^{\text{awake}} \in [0, 1]$ and
expresses a true \emph{local} density, independent of the global node
count $N$ and therefore stable in sparse or non-uniform deployments.
The combination of normalised coordinates and local awake-neighbour
density gives the policy direct access to the spatial geometry that
drives the coverage objective, addressing the limitation of
distance-to-sink as the sole geometric feature. For $N = 50$ nodes,
$|\mathbf{s}^t| = 450$.

\subsection{Factorized Multi-Binary Action Formulation}
\label{sec:action_factor}

The full joint action space $\{0,1\}^N$ has cardinality $2^N$, which for
$N = 50$ exceeds $10^{15}$—well beyond the regime in which a flat-action
DQN or Double DQN over a single discrete output head can be applied
directly. To avoid this combinatorial intractability, we adopt a
\textit{factorized per-node binary scheduling} formulation with
parameter sharing across nodes: a single Q-network outputs two values
per node, corresponding to $Q(s, a_i = 0)$ (SLEEP) and $Q(s, a_i = 1)$
(AWAKE), and the raw scheduling vector is assembled by combining the
per-node greedy decisions:

\begin{equation}
  a_i^t = \arg\max_{a \in \{0, 1\}}\, Q_\theta(\mathbf{s}^t, a)_i,
  \quad i = 1,\ldots,N.
  \label{eq:factorized_argmax}
\end{equation}

The proposed method is therefore a \textit{factorized DDQN scheduler}
rather than a flat-action DDQN over the full joint action space. The
factorization makes the learning problem tractable but introduces two
well-known pitfalls that the methodology addresses explicitly. First,
per-node Q-heads do not jointly optimize correlated decisions; this is
mitigated through reward terms (coverage, balance) that aggregate over
nodes and through the shared backbone, which couples per-node value
estimates via shared representations. Second, the factorized argmax
compounds the standard maximization bias of DQN across $N$ independent
argmax operations, motivating the Double DQN target adopted in
Sec.~\ref{sec:ddqn_arch}.

\subsection{Constraint-Aware Action Execution}
\label{sec:exec_layer}

The raw per-node decisions produced by the agent are not always
physically feasible: a near-empty node must charge before it can resume
sensing, and a depleting node should not be allowed to vacate its
spatial role without a healthy neighbour assuming responsibility.
Rather than expressing these requirements through penalty shaping
alone, we route the agent's output through a deterministic
\textit{constraint-aware execution layer} before the physical state
transition is applied.

\textbf{Executed action.} Let $\mathbf{a}_t \in \{0,1\}^N$ denote the
raw per-node action proposed by the DDQN agent. Because charging
feasibility and coverage continuity may modify infeasible or unsafe
decisions, the environment applies a deterministic transformation
$\tilde{\mathbf{a}}_t = \mathcal{T}(\mathbf{s}_t, \mathbf{a}_t)$. The
transition and reward are generated using the executed action
$\tilde{\mathbf{a}}_t$, and the replay buffer stores
$(\mathbf{s}_t,\, \tilde{\mathbf{a}}_t,\, r_t,\, \mathbf{s}_{t+1},\,
d_t)$. This makes the method robust to invalid-action handling:
training targets are always grounded in the action that was actually
executed in the environment, not the one nominally proposed by the
policy. The transformation $\mathcal{T}$ is composed of two
domain-informed feasibility rules, applied in this order.

\textbf{Charging feasibility rule.} Any node $i$ with
$\text{SoC}_i / E_{\max} < 0.2$ enters the charging state ($c_i = 1$)
and has its action set to SLEEP, recovering
$\Delta\text{SoC} = 0.05 \cdot E_{\max}$ per step. It exits charging
once $\text{SoC}_i / E_{\max} \geq 0.95$. This encodes the hardware
constraint that a near-depleted node must take itself offline to avoid
an irreversible deep-discharge cycle, irrespective of the agent's
nominal preference.

\textbf{Coverage continuity rule.} For any AWAKE node $i$ with
$\text{SoC}_i / E_{\max} \leq 0.5$, the nearest non-dead, non-charging,
currently-sleeping neighbour $j^*$ is forcibly set to AWAKE in
$\tilde{\mathbf{a}}_t$:

\begin{equation}
  j^* = \arg\min_{\substack{j \neq i,\; a_j = 0 \\
                             \neg\text{dead}(j),\;
                             \neg\text{charging}(j)}}
         \|\mathbf{p}_i - \mathbf{p}_j\|_2.
  \label{eq:coop}
\end{equation}

This transfers spatial sensing responsibility to a healthier neighbour
before an energy-depleted node goes offline, maintaining local coverage
continuity without explicit agent intervention. We deliberately frame
this rule as a feasibility constraint encoded in $\mathcal{T}$ rather
than as an environment-level ``override'': the agent learns directly
under the executed-action distribution induced by $\mathcal{T}$, and
the feasibility layer reduces the dimensionality of the scheduling
problem the policy must solve in regions of state space where domain
knowledge already determines the correct action.

\textit{Edge cases.} The coverage continuity rule is realised
deterministically in three respects. \emph{(i)~No eligible
neighbour.} If the constrained $\arg\min$ in Eq.~\eqref{eq:coop} is
empty (every other node is dead, charging, or already AWAKE), the
rule is silently skipped for node $i$ and $\tilde{a}_i$ retains the
value already produced by the charging rule. \emph{(ii)~Multiple
low-SoC nodes targeting the same neighbour.} Low-SoC AWAKE nodes are
processed in ascending order of node index, and any node selected as
$j^*$ within the current step is recorded in a per-step woken-set;
candidates already in this set are excluded from subsequent
$\arg\min$ searches, so each neighbour can be woken at most once per
step. Ties in $\arg\min$ are therefore broken by the index of the
triggering low-SoC node. \emph{(iii)~Neighbour itself near the
charging threshold.} Only nodes that have \emph{already} entered the
charging state (and were therefore added to the charging set by the
charging feasibility rule) are excluded as candidates; a neighbour
whose SoC is close to, but not yet below, the charging threshold
remains eligible and may be woken, after which the next step's
charging rule will activate if its SoC has fallen below the
threshold.

\subsection{Reward Design}
\label{sec:reward}

The scalar reward $r^t$ is a weighted sum of four components:

\begin{equation}
  r^t = w_{\text{cov}} r_{\text{cov}}
      + w_{\text{eng}} r_{\text{eng}}
      + w_{\text{soh}} r_{\text{soh}}
      + w_{\text{bal}} r_{\text{bal}},
  \label{eq:reward}
\end{equation}

with weights
$(w_{\text{cov}}, w_{\text{eng}}, w_{\text{soh}}, w_{\text{bal}})
= (10.0,\, 5.0,\, 1.0,\, 2.0)$.

\textbf{Coverage reward} $r_{\text{cov}}$: The fraction of grid points
covered by alive AWAKE nodes (Section~\ref{sec:network_model}).
Range $[0, 1]$.

\textbf{Energy efficiency reward} $r_{\text{eng}}$: SoC-weighted energy
expenditure, normalised to $[-1, 0]$. Nodes with lower SoC are
penalised more heavily for being active, discouraging depletion of
already-strained batteries:

\begin{equation}
  r_{\text{eng}} = -\,\mathrm{clip}\!\left(
    \frac{\displaystyle\sum_i e_i^t \cdot
          (1 - \text{SoC}_i / E_{\max})}
         {N \cdot e_{\text{awake}} \cdot 2},\;
    0,\; 1\right).
  \label{eq:r_eng}
\end{equation}

\textbf{SoH reward} $r_{\text{soh}}$: Deviation of mean SoH from the
near-perfect threshold of $0.99$, clipped to $[-1, 1]$. Incentivises
preservation of battery health across the fleet:

\begin{equation}
  r_{\text{soh}} = \mathrm{clip}\!\left(
    \overline{\text{SoH}}^{\,t} - 0.99,\; -1,\; 1\right).
  \label{eq:r_soh}
\end{equation}

\textbf{Balance reward} $r_{\text{bal}}$: Negative standard deviation
of SoC fractions across all nodes, clipped to $[-1, 0]$. Penalises
heterogeneous charge distributions, encouraging fair load spreading:

\begin{equation}
  r_{\text{bal}} = \mathrm{clip}\!\left(
    -\sigma_{\text{SoC}}^t,\; -1,\; 0\right).
  \label{eq:r_bal}
\end{equation}

All components are structured as positive-encouraging or
bounded-negative factors to avoid gradient saturation. An additional
penalty of $-10$ is applied at episode termination due to the death
threshold being exceeded.

\subsection{DDQN Learning Architecture}
\label{sec:ddqn_arch}

\textbf{Shared-backbone factorized Q-network.} Both the online (policy)
network $Q_\theta$ and target network $Q_{\theta^-}$ are
parameter-shared MLPs that consume the centralized observation and emit
per-node binary Q-heads:

\begin{equation}
  \mathbb{R}^{9N}
  \xrightarrow{\text{Linear}} 512
  \xrightarrow{\text{ReLU}} 256
  \xrightarrow{\text{ReLU}} \mathbb{R}^{2N}.
  \label{eq:network}
\end{equation}

The output layer produces $2N$ values representing
$Q(\mathbf{s}, a_i{=}0)$ and $Q(\mathbf{s}, a_i{=}1)$ for each node
$i$, reshaped to $(N, 2)$ for per-node argmax. This is the factorized
Q-head referenced in Sec.~\ref{sec:action_factor}—a single shared
backbone is responsible for both representation learning and per-node
value estimation.

\textbf{Experience replay.} Transitions
$(\mathbf{s}_t, \tilde{\mathbf{a}}_t, r_t, \mathbf{s}_{t+1}, d_t)$ are
stored in a circular buffer of capacity $200{,}000$, where
$\tilde{\mathbf{a}}_t$ is the executed action emitted by the
constraint-aware execution layer (Sec.~\ref{sec:exec_layer}).
Mini-batches of size $64$ are sampled uniformly at random. Learning
does not begin until the buffer contains at least $500$ transitions.

\textbf{Double DQN target on the factorized Q-head.} The key
distinction from standard DQN is the decoupling of action selection and
value evaluation, applied independently to each per-node argmax:

\begin{align}
  a_i^* &= \arg\max_{a}\, Q_\theta(\mathbf{s}', a)_i,
           \quad \forall\, i \in [N],
           \label{eq:ddqn_select} \\
  y     &= r + \gamma (1 - d) \cdot
           \frac{1}{N} \sum_{i=1}^{N} Q_{\theta^-}(\mathbf{s}', a_i^*)_i.
           \label{eq:ddqn_target}
\end{align}

In standard DQN (the ablation baseline), the target network both
selects and evaluates the next-state action at every node, so the
per-node maximization bias compounds additively across the $N$ argmax
operations of the factorized head:

\begin{equation}
  y^{\text{DQN}} = r + \gamma (1 - d) \cdot
  \frac{1}{N} \sum_{i=1}^{N} \max_{a}\, Q_{\theta^-}(\mathbf{s}', a)_i.
  \label{eq:dqn_target2}
\end{equation}

\textbf{Loss and optimization.} Mean squared error over the mean
per-node Q-value at the executed action:

\begin{equation}
  \mathcal{L}(\theta) =
  \left(\frac{1}{N}\sum_i Q_\theta(\mathbf{s},\, \tilde{a}_i)_i - y\right)^{\!2}.
  \label{eq:loss}
\end{equation}

Optimized with Adam ($\text{lr} = 1 \times 10^{-4}$). Gradient norms
are clipped to $10.0$ before each parameter update.

\textbf{Target network update.} The target network parameters
$\theta^-$ are hard-copied from $\theta$ every $500$ learning steps.

\textbf{Epsilon-greedy exploration.} During training, the raw per-node
action is replaced with a uniform Bernoulli sample with probability:

\begin{equation}
  \varepsilon(k) = \varepsilon_{\text{end}} +
  (\varepsilon_{\text{start}} - \varepsilon_{\text{end}}) \cdot
  \max\!\left(0,\; 1 - \frac{k}{\varepsilon_{\text{decay}}}\right),
  \label{eq:epsilon}
\end{equation}

with $\varepsilon_{\text{start}} = 1.0$,
$\varepsilon_{\text{end}} = 0.05$,
$\varepsilon_{\text{decay}} = 50{,}000$ learning steps. Evaluation runs
use $\varepsilon = 0$. Exploration noise is applied to the raw action
$\mathbf{a}_t$; the executed action $\tilde{\mathbf{a}}_t$ is always
feasibility-respecting by construction of $\mathcal{T}$. The complete
training procedure is summarized in Algorithm~\ref{alg:ddqn_wsn}.

\begin{algorithm}[!t]
\caption{DDQN-WSN Training Procedure\label{alg:ddqn_wsn}}
\footnotesize
\begin{tabbing}
\hspace{1.6em}\=\hspace{1.6em}\=\hspace{1.6em}\=\hspace{20em}\=\kill
\textbf{Input:} $N$, $T_{\max}$, episodes $E$, buffer capacity,
  batch $b$, update freq.\ $K$, $\gamma$, $\varepsilon$ schedule\\
\textbf{Init:} $Q_\theta$ (online) and $Q_{\theta^-}$ (target)
  with equal random weights;\\
\> $\mathcal{B} \leftarrow \emptyset$;\quad
  $\varepsilon \leftarrow \varepsilon_{\mathrm{start}}$\\[2pt]
\textbf{for} $e = 1$ \textbf{to} $E$ \textbf{do}\\
\> $s_0 \leftarrow \texttt{env.reset()}$
  \`\textit{// randomize positions; reset batteries}\\
\> $t \leftarrow 0$;\quad $\mathit{done} \leftarrow \mathrm{False}$\\
\> \textbf{while not} $\mathit{done}$ \textbf{do}\\
\>\> with prob.\ $\varepsilon$:\quad
     $\mathbf{a}_t \leftarrow \mathrm{random} \in \{0,1\}^N$\\
\>\> otherwise:\quad
     $\mathbf{a}_t \leftarrow \arg\max_{a} Q_\theta(s_t,\cdot)$
     (per-node)\\
\>\> $\tilde{\mathbf{a}}_t \leftarrow
     \mathcal{T}(s_t,\mathbf{a}_t)$
  \`\textit{// constraint-aware execution}\\
\>\> $s_{t+1}, r_t, \mathit{done}
     \leftarrow \texttt{env.step}(\tilde{\mathbf{a}}_t)$
  \`\textit{// transition on executed action}\\
\>\> $\mathcal{B}.\texttt{push}(s_t,\tilde{\mathbf{a}}_t,r_t,s_{t+1},
     \mathit{done})$\\
\>\> \textbf{if} $|\mathcal{B}| \geq \mathit{min\_replay\_size}$
     \textbf{then}\\
\>\>\> $(s,\tilde{\mathbf{a}},r,s',d) \leftarrow
       \mathcal{B}.\texttt{sample}(b)$\\
\>\>\> $a_i^* \leftarrow \arg\max_{a} Q_\theta(s',\cdot)_i$
  \`\textit{// online selects, per-node}\\
\>\>\> $y \leftarrow r + \gamma(1{-}d)\cdot
       \tfrac{1}{N}\sum_i Q_{\theta^-}(s',a_i^*)_i$
  \`\textit{// target evaluates}\\
\>\>\> $\mathcal{L} \leftarrow
       \mathrm{MSE}\!\left(\tfrac{1}{N}\sum_i
       Q_\theta(s,\tilde{a}_i)_i,\;y\right)$\\
\>\>\> $\theta \leftarrow \texttt{Adam}(\nabla_\theta\mathcal{L})$;\quad
       clip $\|\nabla\| \leq 10.0$\\
\>\>\> every $K$ steps:\quad $\theta^- \leftarrow \theta$\\
\>\> \textbf{end if}\\
\>\> $\varepsilon \leftarrow \varepsilon_{\mathrm{end}}
     + (\varepsilon_{\mathrm{start}}{-}\varepsilon_{\mathrm{end}})
     \cdot \max(0,\;1{-}\mathit{steps}/\varepsilon_{\mathrm{decay}})$\\
\>\> $t \leftarrow t + 1$\\
\> \textbf{end while}\\
\> Log: reward, coverage,
  $\overline{\mathrm{SoH}}$, alive fraction\\
\textbf{end for}
\end{tabbing}
\end{algorithm}

% -----------------------------------------------------------------------
\section{Experimental Setup}
\label{sec:experiments}
% -----------------------------------------------------------------------

\subsection{Network and Training Configuration}

All experiments use the configuration listed in
Table~\ref{tab:config}. The same seed controls node position
initialization, neural network weight initialization, and exploration
noise, ensuring full reproducibility.

\begin{table}[!t]
\caption{Experimental Configuration\label{tab:config}}
\centering
\footnotesize
\setlength{\tabcolsep}{4pt}
\begin{tabular}{lc}
\toprule
\textbf{Parameter} & \textbf{Value} \\
\midrule
Number of nodes $N$                    & 50 \\
Arena size                             & $500\,\text{m} \times 500\,\text{m}$ \\
Sink position                          & $(250,\,250)$ \\
Sensing radius $r_s$                   & $100\,\text{m}$ \\
Coverage grid resolution               & $20 \times 20$ points \\
Max steps per episode $T_{\max}$       & 1000 \\
Death threshold $\delta$               & 0.3 \\
Battery capacity $E_{\max}$            & 100 \\
Awake energy draw $e_{\text{awake}}$   & 1.0 \\
Sleep energy draw $e_{\text{sleep}}$   & 0.01 \\
Cycle degradation $k_{\text{cycle}}$   & $5 \times 10^{-5}$ \\
DoD exponent $\alpha$                  & 1.2 \\
Calendar decay $k_{\text{cal}}$        & $5 \times 10^{-7}$ \\
Charging threshold                     & 0.20 (SoC fraction) \\
Charging rate                          & 0.05 (SoC fraction/step) \\
Cooperative wake-up SoC threshold      & 0.50 \\
Training episodes                      & 500 \\
Batch size $b$                         & 64 \\
Learning rate                          & $1 \times 10^{-4}$ \\
Discount factor $\gamma$               & 0.99 \\
Replay buffer capacity                 & $200{,}000$ \\
Min replay size                        & 500 \\
Target update frequency $K$            & 500 steps \\
$\varepsilon_{\text{start}}$ / $\varepsilon_{\text{end}}$ / $\varepsilon_{\text{decay}}$
                                       & 1.0 / 0.05 / $50{,}000$ \\
Q-Network hidden layers                & [512,~256] \\
Optimiser                              & Adam \\
Gradient clip norm                     & 10.0 \\
Random seed                            & 42 \\
\bottomrule
\end{tabular}
\end{table}

\subsection{Compared Methods}

Two agents are evaluated under identical environmental conditions:

\begin{itemize}
  \item \textbf{DDQN-WSN} (proposed): Double DQN with decoupled action
    selection/evaluation (Section~\ref{sec:ddqn_arch}). Action
    selection uses the online network; Q-value estimation uses the
    target network.
  \item \textbf{DQN} (ablation): Standard Deep Q-Network. Action
    selection and evaluation are both performed by the target network via
    direct maximization. Identical architecture, hyperparameters, and
    environment.
\end{itemize}

Both agents are trained from scratch with seed 42. Evaluation runs use
$\varepsilon = 0$ (fully greedy) to measure the learned policy without
exploration noise.

\subsection{Evaluation Protocol (Training-Time)}
\label{sec:eval_protocol}

The empirical analysis in this paper characterises the converged
behaviour observed during training. Concretely, the converged policy
of each agent is represented by the moving average of the final ten
training episodes---a standard proxy for converged policy quality in
episodic settings. Metrics reported include: coverage fraction, mean
SoH, alive fraction, episode reward, and within-episode network
lifetime, defined operationally as
\begin{equation}
  T_{\text{life}} = \min\!\big\{ t \in [1, T_{\max}] :
  \text{dead\_fraction}(t) > \delta \big\},
  \label{eq:lifetime}
\end{equation}
capped at $T_{\max} = 1000$ steps when the death threshold is never
breached within the episode. Because per-node SoH is strictly
non-increasing, the within-episode alive fraction is monotone, so the
event $\text{alive\_fraction}_{\text{end}}(e) \geq 1 - \delta$
recorded at the end of every training episode $e$ implies
$T_{\text{life}}(e) = T_{\max}$ for that episode. Training curves
smoothed with a 10-episode moving average are reported to assess
convergence stability across the full 500-episode training run.

\textit{Held-out evaluation.} An external benchmark of the converged
policy should additionally report performance on a separate
held-out evaluation set: the trained policy is frozen, exploration
is disabled ($\varepsilon = 0$), and a number of deterministic test
episodes are rolled out on previously unseen random topologies. A
reasonable minimum protocol is five independent seeds with twenty
deterministic test episodes per seed (with means and standard
deviations, ideally $95\%$ confidence intervals); ten seeds with
fifty test episodes each yields tighter confidence intervals when
compute permits. Such a held-out evaluation constitutes the
recommended protocol for an external benchmark and is planned as part
of the deployment validation. The training-time results presented in
this paper therefore characterise convergence behaviour rather than
generalisation, and should be read in that scope.

% -----------------------------------------------------------------------
\section{Results and Discussion}

\subsection{Overall Performance Comparison}

\emph{All numerical results in this section are training-time moving
averages computed over the final ten episodes of each agent's
500-episode training run
(Sec.~\ref{sec:eval_protocol}); they characterise convergence
behaviour rather than held-out generalisation, and a separate
deterministic held-out evaluation is planned as part of deployment
validation.}

Table~\ref{tab:dqn_vs_ddqn_wsn} summarises the final performance of
DDQN-WSN versus standard DQN, represented by the moving average of
the final 10 training episodes of each agent's converged policy.
Mean SoC at termination is derived from the final entry of the
per-episode \texttt{mean\_soc} series recorded during training.

\begin{table}[h]
\centering
\small
\setlength{\tabcolsep}{4pt} % reduce column spacing

\begin{tabularx}{\linewidth}{|X|c|c|X|}
\hline
\textbf{Metric} & \textbf{DQN} & \textbf{DDQN-WSN} & \textbf{Difference} \\
\hline

Mean Episode Reward & 8297.79 & 8305.82 & +8.03 (+0.10\%) \\
\hline

Final Coverage (\%) & 96.45\% & 96.36\% & $-0.09$ pp \\
\hline

Mean SoH at termination & 0.99968 & 0.99968 & $\approx$ Parity \\
\hline

Network Lifetime $T_{\text{life}}$ (steps) & $\geq 1000$ ($T_{\max}$) & $\geq 1000$ ($T_{\max}$) & Parity \\
\hline

\end{tabularx}

\caption{Comparison between DQN and DDQN-WSN performance metrics}
\label{tab:dqn_vs_ddqn_wsn}
\end{table}

\subsection{Per-Metric Analysis}

\subsubsection{Energy Consumption}

Per-episode energy consumption is tracked as the mean SoC drain across
the network, $\Delta\overline{\mathrm{SoC}} = \overline{\mathrm{SoC}}_0
- \overline{\mathrm{SoC}}_{T}$, where $\overline{\mathrm{SoC}}_0$ and
$\overline{\mathrm{SoC}}_{T}$ are the mean state-of-charge fractions
at episode start and end, respectively. Averaged over the final ten
training episodes, DDQN-WSN exhibits a mean drain of $\approx 0.418$
and standard DQN $\approx 0.413$ (Fig.~\ref{fig:energy_consumption}),
i.e.\ both agents discharge roughly $41$--$42\%$ of fleet-wide SoC
over the 1000-step episode. This is consistent with the
$\sim\!24$--$25$ awake nodes per step implied by the coverage analysis
below: with $e_{\text{awake}} = 1.0$ per awake-node-step and
$E_{\max} = 100$, an awake population of $\sim\!25$ nodes drains
roughly $25 \cdot 1000 / (50 \cdot 100) = 0.50$ of fleet SoC, of which
a portion is recovered by the charging state machine.

The SoC-weighted energy component of the reward
function~(Eq.~\eqref{eq:r_eng}) penalises activity on nodes that are
already low on charge, so the agent is incentivised to route sensing
responsibility toward higher-SoC nodes whenever the constraint layer
permits. The near-identical drain profiles of the two agents reflect
the parity observed across the other converged metrics
(Table~\ref{tab:dqn_vs_ddqn_wsn}); a finer-grained per-node and
per-step energy breakdown would require extending the trainer's
logging schema and is left for future work.


\begin{figure}[h]
\centering
\includegraphics[width=\linewidth]{energy_consumption (1).png}
\caption{Per-episode energy consumption, defined as the mean SoC drain
across the network within an episode
$\Delta\overline{\mathrm{SoC}} = \overline{\mathrm{SoC}}_0 -
\overline{\mathrm{SoC}}_{T}$ (dimensionless, $[0,1]$). DDQN-WSN vs.\
standard DQN over 500 training episodes; curves are smoothed with a
10-episode moving average.}
\label{fig:energy_consumption}
\end{figure}

\subsubsection{Network Coverage}

Both agents converge to high and stable coverage fractions over the 500-episode training run, with DDQN-WSN and DQN achieving final coverage of 96.36\% and 96.45\% respectively. The cooperative wake-up mechanism plays a key role in maintaining spatial coverage continuity: as individual nodes approach the charging threshold and are forced to SLEEP, their nearest available neighbour is automatically activated, preventing localised coverage gaps without requiring explicit agent intervention.

Coverage is computed over a fixed $20 \times 20$ grid (400 sample points) with a sensing radius of 100 m. A single fully-charged node at the arena centre can cover approximately 12.57\% of the arena ($\pi \times 100^2 / (500 \times 500) \approx 0.1257$). Maintaining approximately 96.4\% coverage with 50 uniformly deployed nodes requires approximately 24 to 25 awake nodes at any given step, assuming optimal spatial distribution — solving $(1 - 0.1257)^n = (1 - 0.964)$ for $n$ yields $n \approx 24.8$.



\begin{figure}[h]
\centering
\includegraphics[width=\linewidth]{coverage.png}
\caption{Per-Episode Coverage Fraction — DDQN-WSN vs. DQN (500 training episodes)}
\label{fig:Per-Episode Coverage Fraction — DDQN-WSN vs. DQN (500 training episodes}
\end{figure}

\subsubsection{Network Lifetime}

Network lifetime is reported as the within-episode quantity
$T_{\text{life}}$ defined in Eq.~\eqref{eq:lifetime}: the first
timestep at which more than $\delta N = 15$ of the $50$ nodes have
permanently failed, or $T_{\max} = 1000$ if no such timestep occurs.
This distinguishes operational survival time inside a single
deployment (measured in steps) from the training horizon (measured
in episodes).

Across the full 500-episode training run, both DDQN-WSN and standard
DQN reach the end of every episode with at least $1 - \delta = 70\%$
of nodes alive. Since per-node SoH degradation is irreversible and
the within-episode alive fraction is therefore monotone
non-increasing, the absence of an end-of-episode breach is sufficient
to conclude that the death threshold was not crossed at any earlier
step either: $T_{\text{life}} = T_{\max} = 1000$ steps for every
training episode of both agents. The SoH-aware degradation model
ensures that nodes cannot be repeatedly deep-discharged without
long-term health consequences, so this within-episode lifetime is a
more demanding metric than a simple SoC-depletion count would be.

\subsubsection{SoC Distribution and Load Balance}



The balance reward $r_{\text{bal}} = -\sigma_{\text{SoC}}$ directly penalises unequal charge distribution. An optimal policy would maintain all nodes at approximately equal SoC — maximising the homogeneity of the residual energy landscape and extending the point at which any single node hits the death threshold.

\subsubsection{Alive Fraction Over Time}



Unlike SoC, which can be partially recovered through the charging state machine, SoH degradation is irreversible. The two mechanisms combine to produce two failure modes: short-term SoC exhaustion (recoverable if charging kicks in) and long-term SoH collapse (permanent). DDQN-WSN's health-preserving reward component specifically targets the second failure mode.

\subsubsection{Battery SoH Preservation}


\begin{figure}[h]
\centering
\includegraphics[width=\linewidth]{battery_health.png}
\caption{Mean Battery State of Health — DDQN-WSN vs. DQN (500 training episodes)}
\label{fig:battery_health}
\end{figure}

As shown in Figure 4, both agents maintain near-perfect battery health throughout the entire 500-episode training run, with mean SoH remaining flat at approximately 0.9997 from episode 50 onward. The two curves are indistinguishable, confirming that both DDQN-WSN and DQN learn to avoid deep-discharge patterns that would trigger meaningful SoH degradation. Final mean SoH at episode 500: DDQN-WSN 0.99968, DQN 0.99968 — complete parity.

This outcome is a direct consequence of the super-linear DoD degradation model: the $r_{\text{soh}}$ reward component, combined with the $\alpha = 1.2$ exponent, creates a strong incentive for both agents to adopt shallow-discharge scheduling patterns early in training and sustain them throughout.

The cycle degradation model with $\alpha = 1.2 > 1$ creates a super-linear penalty for deep discharges. An agent that repeatedly deep-discharges a node (SoC: 100 $\rightarrow$ 0) accumulates $k_{\text{cycle}} \cdot 1.0^{1.2} = 5 \times 10^{-5}$ health loss per cycle, versus $k_{\text{cycle}} \cdot 0.1^{1.2} \approx 3.8 \times 10^{-6}$ for a shallow 10\% discharge — a 13$\times$ difference. A learned policy that distributes load and avoids deep discharges can therefore extend SoH significantly relative to a greedy policy.

\subsection{Ablation Study}

\subsubsection{Double DQN vs. Standard DQN}

The core algorithmic difference between DDQN-WSN and the DQN baseline is the target computation.

\[
y^{\text{DQN}} = r + \gamma \cdot \frac{1}{N} \sum_i \max_a Q_{\theta^-}(s', a)_i
\]

This introduces a positive bias:
\[
\mathbb{E}[\max_a Q_{\theta^-}(s', a)] \geq \max_a \mathbb{E}[Q_{\theta^-}(s', a)]
\]

DDQN corrects this by using the online network for selection and the target network for evaluation only:

\[
y^{\text{DDQN}} = r + \gamma \cdot \frac{1}{N} \sum_i Q_{\theta^-}(s', a_i^*)_i
\]

\[
a_i^* = \arg\max_a Q_\theta(s', a)_i
\]


\begin{figure}[h]
\centering
\includegraphics[width=\linewidth]{throughput.png}
\caption{Network Throughput (Coverage Proxy) — DDQN-WSN vs. DQN (500 training episodes)}
\label{fig:Network Throughput (Coverage Proxy) — DDQN-WSN vs. DQN (500 training episodes}
\end{figure}

Both agents converge to comparable throughput levels, with the training curves largely overlapping after episode $\sim$50. The DDQN-WSN curve exhibits marginally lower variance in the mid-training range (episodes 100–300). Final mean episode reward: DDQN-WSN 8,305.82 vs. DQN 8,297.79 — a difference of +8.03 (+0.10\%), reflecting effective parity at convergence.

\subsubsection{Constraint-Layer Analysis: RL Contribution vs.\ Heuristic Override}

A natural question is whether the converged performance of DDQN-WSN
is driven primarily by the learned policy or by the deterministic
constraint-aware execution layer~$\mathcal{T}$ alone. The
constraint layer enforces only two local feasibility rules---a
charging rule that takes near-depleted nodes offline, and the
coverage-continuity rule of Eq.~\eqref{eq:coop} that promotes a
single nearest neighbour---neither of which contains any explicit
notion of global coverage, energy efficiency, or charge-balance
fairness. Three variants are therefore of structural interest:

\begin{itemize}
  \item \textbf{Constraint-Only.} Raw actions are sampled uniformly
    at random from $\{0,1\}^N$ and filtered by~$\mathcal{T}$. The
    charging rule prevents the most catastrophic deep-discharge
    deaths and the cooperative wake-up rule maintains a local
    handover when an AWAKE node depletes; however, the random
    actions provide no global spatial coordination, so coverage is
    expected to fluctuate well below the $96\%+$ level attained by
    DDQN-WSN.
  \item \textbf{Unconstrained DDQN.} The DDQN policy is trained and
    executed without~$\mathcal{T}$, so raw actions are applied
    directly. The agent must learn the survivability constraints
    from reward signals alone; in regions of state space where a
    single misstep triggers an irreversible deep-discharge cycle,
    boundary-case starvation events are expected to increase
    relative to the full system.
  \item \textbf{Full DDQN-WSN.} The combination of the learned
    policy and~$\mathcal{T}$ is expected to outperform both
    ablations: $\mathcal{T}$ removes the worst feasibility failures
    so that the policy can specialise on global spatial
    orchestration and long-horizon SoH preservation, while the
    learned policy supplies the global coordination that the
    deterministic rules lack.
\end{itemize}

\textit{Planned ablation (Table TBD).} A full quantitative ablation
table reporting coverage, end-of-episode failure rate, mean reward,
mean SoH, and within-episode lifetime for each of the three variants
above is part of the deployment-validation phase and is left for a
companion experimental study. The qualitative argument given here
should not be read as a substitute for those numbers.
\subsection{Discussion: Role of the Double DQN Target}

The converged-policy results in Table~\ref{tab:dqn_vs_ddqn_wsn} show
that DDQN-WSN and standard DQN reach effective parity under the
present 500-episode setting; the rationale for adopting the Double
DQN target on the factorized Q-head is therefore best read as a
structural argument rather than a claim of empirical superiority at
this horizon. Three interacting factors are relevant:

\textbf{Theoretically lower target-estimation bias:} Decoupling action
selection (online network) from value evaluation (target network)
removes the per-node maximization bias that would otherwise compound
additively across the $N = 50$ argmax operations of the factorized
Q-head (cf.\ Eqs.~\eqref{eq:ddqn_target} and \eqref{eq:dqn_target2}).
This property is not directly measured here but is a well-established
benefit of the Double DQN target in higher-variance or
longer-horizon settings.

\textbf{Domain-informed reward shaping:} The four-component reward
function ensures that no single objective dominates. The high
coverage weight ($w_{\text{cov}} = 10$) provides the primary signal,
while the energy and balance components act as regularisers that
discourage degenerate policies (e.g.\ always-awake or round-robin
schedules with no spatial awareness). The SoH component, even at low
weight ($w_{\text{soh}} = 1$), aligns the policy with shallow-discharge
patterns under the SoC/SoH-aware degradation model.

\textbf{Environment-level cooperation:} The cooperative wake-up rule in
the constraint-aware execution layer transfers sensing responsibility
to a healthy neighbour whenever an AWAKE node falls below the
cooperative-wake SoC threshold, reducing the dimensionality of the
scheduling problem that the agent must solve in the
near-depletion regime.

\subsection{Limitations and Broader Impact}

\textbf{Centralized control:} The current formulation assumes a centralized scheduler with full observability of all 50 nodes simultaneously. In large-scale deployments, communication overhead for state aggregation may be prohibitive. Decentralized and partially observable extensions are a natural direction for future work.

\textbf{Higher-order joint dependency modeling:} Our factorized DDQN approach can be used to overcome the intractability of the joint action space by estimating Q-values per node, but it does not make it possible to explicitly model node interaction dependencies. This model is based on the approximation of the joint scheduling at the output layer, instead of explicitly modeling the conditional relationships between neighboring actions. These higher-order joint dependencies can be captured by implementing more sophisticated multi-agent coordination architectures like Branching DQN, value factorization networks (e.g., QMIX), or explicit decentralized MARL communication protocols---without just relying on cooperative wake-up heuristics in the environment.

\textbf{Static topology:} Node positions are randomized at episode start but remain fixed within an episode. Mobile sensor networks and environments with node addition/removal are not currently supported.

\textbf{Homogeneous hardware:} All nodes share identical battery parameters ($E_{\max}$, $k_{\text{cycle}}$, $\alpha$). Real deployments may involve heterogeneous hardware with differing capacity and degradation characteristics.

\textbf{Simulation fidelity:} The battery model captures the dominant degradation mechanisms (DoD cycling and calendar aging) but omits temperature effects, manufacturing variance, and charge-rate-dependent aging. Bridging the simulation-to-real gap for physical deployment would require calibration against real cell discharge data.



% -----------------------------------------------------------------------
\section{Conclusion}
\label{sec:conclusion}
% -----------------------------------------------------------------------

This paper presented DDQN-WSN, a constraint-aware factorized Double
Deep Q-Network agent for scheduling sensor node sleep/wake cycles in a
50-node wireless sensor network. The system is built on a surrogate,
empirical SoC/SoH-aware battery degradation model tracking both State
of Charge (SoC) and irreversible State of Health (SoH) loss via
cycle-based and calendar aging. To avoid the combinatorial
intractability of the full $2^N$ joint action space, scheduling is
formulated as $N$ parameter-shared per-node binary decisions, and the
agent's raw action is routed through a deterministic constraint-aware
execution layer $\mathcal{T}$ that enforces charging feasibility and
coverage continuity; the resulting executed action
$\tilde{\mathbf{a}}_t$ drives both the environment transition and the
contents of the replay buffer, eliminating the invalid-action ambiguity
that arises when feasibility rules are applied as post-hoc overrides.

The DDQN agent observes a nine-feature per-node state vector—including
normalized coordinates and a local awake-neighbour density score that
exposes spatial geometry directly to the policy—and optimizes a
four-component reward balancing area coverage, SoC-weighted energy
efficiency, battery health preservation, and charge-level fairness.
Direct comparison against standard DQN under identical conditions
shows that, within the present 500-episode training horizon,
DDQN-WSN achieves performance comparable to DQN on every converged
metric ($96.36\%$ vs.\ $96.45\%$ coverage, $0.99968$ mean SoH, no
within-episode network collapse observed across all 500 training
episodes), while providing a more robust target-estimation mechanism
for the factorized per-node scheduling problem: by selecting the
next-state action with the online network and evaluating it with the
target network, the per-node maximization bias does not compound
additively across the $N$ argmax operations of the factorized
Q-head. The empirical benefit of this property is expected to be most
visible under longer-horizon or higher-variance deployment settings
than those exercised here.

Several architectural directions extend the present centralized
factorized scheduler and are planned for future work:

(i)~\textbf{Graph Neural Networks (GNNs) and Graph-Aware Policy Learning}, using a GNN backbone to capture the dynamics of the spatial topology, multi-hop communication links and coverage overlap naturally.

(ii)~\textbf{Decentralized Multi-Agent RL (MARL)}: breaking the centralized scheduling algorithm into explicit per-node agents, which communicate with one another locally to make scheduling decisions, without a single point of failure and scalable to massive networks.

(iii)~\textbf{Proximal Policy Optimization (PPO) and Actor-Critic Methods}, moving beyond the constrained value-based DDQN scheduling methods to more powerful policy gradient methods that can afford continuous actions (e.g., tunable transmission power) and more complex stochastic policies.

(iv)~\textbf{Advanced Sim-to-Real Battery Models}, which incorporates current empirical surrogate degradation models with high fidelity electrochemical dynamic models, non-linear thermal coupling, and real-world hardware calibration.

% -----------------------------------------------------------------------
\begin{thebibliography}{15}
\bibliographystyle{IEEEtran}

\bibitem{mohammed2024}
M.~Mohammed, R.~Khalil, Z.~O.~Najah, and Z.~D.~Najah,
``Federated learning-based clustering for energy efficiency in wireless
sensor networks,''
2024.

\bibitem{burger2023}
M.~Burger,
``Soft K-Means clustering for energy-efficient wireless sensor
networks,''
2023.

\bibitem{nguyen2024}
N.~Duy~Tri and T.~T.~H.~Le,
``Energy-efficient cluster routing in wireless sensor networks using
fuzzy logic,''
2024.

\bibitem{kaur2021}
H.~Kaur and U.~Kaur,
``Modified LEACH protocol for energy efficiency in wireless sensor
networks,''
2021.

\bibitem{xue2020}
X.~Lin,
``Deep Q-network-based node management for wireless sensor networks,''
2020.

\bibitem{li2023}
S.~Li and Z.~Zhenyu,
``QDNC packet scheduling for wireless sensor networks,''
2023.

\bibitem{muhammad2023}
A.~Muhammad,
``Reinforcement learning-based sleep/wake scheduling in wireless sensor
networks,''
2023.

\bibitem{chang2023}
C.~Lee,
``Particle swarm optimization with fuzzy clustering for WSN energy
optimization,''
2023.

\bibitem{su2023}
S.~Su,
``Q-learning-based routing protocol for wireless sensor networks,''
2023.

\bibitem{wang2023}
X.~Wang, H.~Chen, and S.~Li,
``A reinforcement learning-based sleep scheduling algorithm for
compressive data gathering in wireless sensor networks,''
\emph{EURASIP Journal on Wireless Communications and Networking},
vol.~2023, no.~28, 2023.

\bibitem{Jin2022}
W.~Jin, J.~Sun, K.~Chi, and S.~Zhang,
``Deep reinforcement learning based scheduling for minimizing age of
information in wireless powered sensor networks,''
\emph{Computer Communications}, vol.~191, pp.~1--10, 2022.

\bibitem{nishant2025}
T.~A.~H.~Nishat, J.-H.~Jeong, H.~Jo, S.~Xia, and J.~Liu,
``Reinforcement learning for adaptive battery management of structural
health monitoring IoT sensor network,''
\emph{Applied Energy}, vol.~390, p.~125731, 2025.

\bibitem{Manogaran2025}
N.~Manogaran, M.~T.~M.~Raphael, R.~Raja, A.~K.~Jayakumar,
M.~Nandagopal, B.~Balusamy, and G.~Ghinea,
``Developing a novel adaptive double deep Q-learning-based routing
strategy for IoT-based wireless sensor network with federated
learning,''
\emph{Sensors}, vol.~25, no.~10, p.~3084, 2025.

\bibitem{Kim2024}
H.~Kim, Y.-J.~Kim, and W.-T.~Kim,
``Deep reinforcement learning-based adaptive scheduling for wireless
time-sensitive networking,''
\emph{Sensors}, vol.~24, no.~16, p.~5281, 2024.

\bibitem{Jaber2024}
M.~M.~Jaber \emph{et al.},
``Q-learning based task scheduling and energy-saving MAC protocol for
wireless sensor networks,''
\emph{Wireless Networks}, vol.~30, pp.~4989--5005, 2024.

\end{thebibliography}

\vfill

% -----------------------------------------------------------------------
% Author biographies
% -----------------------------------------------------------------------



\end{document}